\par\SourceLeaf{148}\noindent
Aus den frühern Erfahrungen jedoch [S. 139 unten] hätte man, unter Berücksichtigung des Umstandes, dass in den ersten Decennien die Gesellschaft einen viel kleinern Umfang hatte als 1794, und dass jene Erfahrungen einer Zeit angehörten, wo der entsprechende Beharrungszustand als noch nicht ganz erreicht angesehen werden muss, schliessen können, dass man fortan eine die Durchschnittszahl 13 weit überschreitende Witwenzahl für sehr wohl möglich halten müsse. (Dass bei allem, was ich hier gesagt habe, die nach \textsc{Kritter}’s Gutachten noch erforderliche Vergrösserung um \(\tfrac16\), wegen der Waisenpensionen, noch nicht mit einbegriffen ist, wird man nicht übersehen dürfen).

Die \textsc{Kritter}’sche S. [140] Nr. 2 gerügte Behauptung ist zwar durch die Erfahrung genugsam widerlegt: es ist jedoch nicht überflüssig, zu der eigentlichen etwas versteckt liegenden Quelle des Irrthums hinaufzusteigen. Bei aller Anwendung des Calculs sowohl auf Gegenstände der Natur als auf sociale Verhältnisse, pflegen die Erfahrungsdata selten in der reinen Gestalt, wie man sie eigentlich braucht, aufzutreten, sondern fast immer mehr oder weniger behaftet mit Störungen oder Schwankungen, die in ihrem Wechsel keiner Regel gehorchen, und man sucht dann, wie jedermann weiss, den daraus entstehenden Nachtheil wenn auch nicht aufzuheben, doch so viel thunlich zu vermindern, dass man aus vielen einzelnen Resultaten das Mittel nimmt. Man rechnet darauf, dass bei einer solchen Benutzung einer grossen Zahl von Fällen die zufälligen Schwankungen einander grösstentheils compensiren, und legt dann dem Mittelwerthe eine desto grössere Zuverlässigkeit bei, je mehr partielle Resultate zugezogen sind. Dieses ist auch im allgemeinen vollkommen richtig, und durch consequente weitere Entwicklung und umsichtige Ausbeutung dieses Princips sind besonders in den Naturwissenschaften nicht selten die belohnendsten Früchte, selbst glanzende Resultate, gewonnen. Allein die Sicherheit des Grundprincips beruhet auf einer wesentlichen Bedingung, die, häufig genug, auch von Gelehrten vom Fach ausser Acht gelassen wird, und die darin besteht, dass die an den einzelnen Beobachtungen oder Erfahrungen haftenden regellosen Störungen oder Schwankungen von einander ganz unabhängig sein müssen. Das Urtheil, ob eine solche Unabhängigkeit vorhanden sei oder nicht, kann zuweilen sehr schwierig und ohne tiefes Eindringen in das Sachverhältniss unmöglich sein, und wenn darüber Zweifel zurückbleiben, so wird auch das den Endresultaten beizulegende Gewicht ein precäres sein.

Wäre z. B. die Rede von einem meteorologischen Elemente etwa von der Menge des an einem bestimmten Orte jährlich fallenden Regens, so ist diese bekanntlich in verschiedenen Jahren sehr ungleich; der durch die allgemeinen örtlichen Verhältnisse des Platzes bedingte Normalwerthwird aber an einem Durchschnitt von zehn Jahren mit viel grösserer Sicherheit erkannt, als wenn man sich bloss an ein einzelnes Jahr halten wollte. Der Grund ist aber der, weil zwischen den in den einzelnen Jahren vorkommenden Abweichungen von dem Normalwerthe kein besonderer Zusammenhang ist, vielmehr, wie auch die Erfahrung bestätigt, eine grosse Minus-Abweichung eben so leicht in einem Jahre vorkommen kann, welches unmittelbar auf ein Jahr mit grosser Plus-Abweichung folgt, wie in jedem andern.

Allein jene wesentliche Bedingung fehlt bei den gezählten Witwen aus auf einander folgenden Jahren, eben weil der Uebergang von einer Zahl zu einer bedeutend verschiedenen nur allmählich geschehen kann. Wenn z. B. in der oben zur Erläuterung angeführten grössern Gesellschaft, wo der \textit{durchschnittliche} jährliche Zugang zu 69 angenommen ist, und eben so gross, nach erreichtem Beharrungszustande, der jährliche Abgang, der Bestand einmal auf 1240 heruntergekommen ist, oder dermalen die negative Abweichung \(-60\) Statt findet, so ist die grösste an Unmöglichkeit grenzende Unwahrscheinlichkeit da, dass im Jahre darauf eine positive Abweichung vom Normalwerthe Statt haben werde. Bei einer kleinen Gesellschaft wie die unsrige sind sehr oft die gezählten Witwen des folgenden Jahres noch ganz die nämlichen wie im vorangegangenen, und selbst nach 10 Jahren wird in der Regel nur der kleinere Theil erneuert sein. Eine Durchschnittszahl aus 10 auf einander folgenden Jahren ist daher noch kein Mittel aus 10 von einander ganz unab%
